\documentclass[a4paper]{moderncv}

% 风格
\moderncvcolor{cerulean}
\moderncvstyle{contemporary}

% 编码（推荐用 xelatex 编译）
\usepackage{xeCJK}
% \setCJKmainfont{PingFang SC}

\usepackage[scale=0.85]{geometry}
\usepackage{setspace}
\setstretch{1.2}

% 个人信息
\name{居}{振梁}
\title{青少年科技辅导员 | 课程体系设计师}
\phone[mobile]{+86-186-0123-9387}
\email{juzhenliang@hotmail.com}
\social[github]{wargrey}
\extrainfo{微信: shadowargrey | 公众号: plteenfun}

\begin{document}

\makecvtitle

%--------------------------------------------------
\section{个人简介}

计算机背景\texttt{STEM}教育者, 
关注全学段跨学科课程设计与信息学奥赛教学。
课程体系立足学生已有知识网络,
自顶向下构建“可见结构 → 可操作结构 $\rightarrow$ 抽象结构”认知路径, 
并以自研教学工具与课程系统支撑规模化落地,
教学经验覆盖 2--11 年级多种教学场景。

%--------------------------------------------------
\section{教学与课程经验}

\cventry{2024--2026}{独立教师}{义乌市青舞艺术培训部有限公司}{金华}{}{
\begin{itemize}
\item 与义乌市福田第二小学合作, 承担三年级信息技术教学及\texttt{Python}社团课教学
\item 构建信息学奥赛教学路径(已有知识 → 可见结构 → 可操作结构 → 竞赛问题),
        教授 4--8 年级\texttt{CSP/J}, 通过精心设计的学生手册引导学生纸笔推演,
        逐步从依赖背模版过渡到独立建模
\item 实验性开设高中段 STEM 暑假体验班, 验证课程在高年级阶段的适配性
\item 将自研教学用游戏引擎写成学位论文
\end{itemize}
}

\cventry{2021--2024}{教务老师}{镇江空白文化传媒有限公司(洛克白STEM)}{镇江}{}{
\begin{itemize}
\item 与镇江培文实验学校合作, 3--6 年级义务教育科学主教老师, 成绩由 30--40 分提升到 70--80 分;
        负责江苏省实验操作大赛第一届和第二届的组织和培训, 金牌、银牌数量总计5个
\item 设计并实施\texttt{STEM}课程体系和跨学科编程教育体系, 
        覆盖小实验、实验探究及综合项目制课程, 并衔接蓝桥青少和信息素养大赛
\item 研发教学用游戏引擎及外围工具链, 系统性降低了\textbf{无基础}学生对\texttt{C++}的恐惧;
        五年级学生从抗拒 C++ 到独立完成《三原色打靶游戏》(跨数学、物理学、语言学三科)仅需16次课
\end{itemize}
}

%--------------------------------------------------
\section{科研与项目经历}

\cventry{2017--2020}{软件工程师}{镇江明润信息科技有限公司}{镇江}{}{
\begin{itemize}
\item 以“清晰优先于形式”的低认知负荷设计原则, 将碎片化的行业经验和知识落地软件：
        从零构建基于\texttt{UWP}平台的疏浚控制系统, 
        以优雅的设计替代了老系统, 随后演化成了我的教学系统
\item 参与船载环境累计 6 个月实地测试, 在复杂工况下验证系统稳定性, 并基于反馈持续迭代优化
\end{itemize}
}

\cventry{2014--Present}{开源贡献者}{Racket 社区}{远程}{}{
\begin{itemize}
\item 构建面向编程教育的知识表达与学习系统, 持续演进超过 10 年；
        包括可视化语义建模\texttt{DSL}、
        初学者友好的辅助工具链(\texttt{C++}构建系统、测试与题库测评系统等)、
        用于编写教材、生成课件素材的写作系统等
\item 开发多项基础设施组件, 涵盖数据压缩(\texttt{zlib}）、安全通信（\texttt{SSH}/\texttt{ASN.1})等方向
\end{itemize}
}

%--------------------------------------------------
\section{培训与证书}

\cventry{2021}{青少年科技辅导员认证}{CACSI}{}{}{}
\cventry{备考}{高中数学教师资格证}{}{}{}{}

%--------------------------------------------------
\section{教育背景}

\cventry{2020--2025}{南京航空航天大学}{计算机科学与技术}{本科}{南京}{}
\cventry{2006--2009}{苏州工学院(原常熟理工学院)}{J2EE企业级应用开发}{大专}{常熟}{}

\end{document}
